\section{Планковский мир}
\label{sec:planck-floor}
\label{ch:carrier-object}
\label{sec:carrier-object}

Масштаб длины, на котором одновременно видны
квантовость, локальная кривизна и причинность, уже записан: $\lengthPlanck$ в \eqref{eq:bh-entropy}.
Вместе с ним время~$\timePlanck=\lengthPlanck/\speedC$ и масса~$m_P=\sqrt{\planckHbar \speedC/\constGrav}$,
а энергия того же масштаба --- \eqref{eq:planck-energy}.
Физика по-прежнему не зависит от человеческих единиц:
важны безразмерные отношения и те же симметрии, что названы выше.

Требование ограниченной плотности уже названо.
Если в каждой точке $x\in\spaceRthree$ сидит своё комплексное число,
в любом конечном причинном объёме конфигураций несчётно много,
а требование уже ограничивает это число.
Место, на котором это выполнено, --- счётное множество узлов.

\begin{definition}[Носитель как множество узлов]
\label{def:carrier-sites}
На микроуровне~$\layerMicro$ \textbf{носителем} имеют в виду счётное множество узлов~$\carrierLattice$.
Время, согласованное с одним шагом сигнала, отождествляют с дискретной шкалой $t\in\spaceZ$.
\end{definition}

Требование конечной скорости уже названо.
На счётном множестве узлов конечное расстояние за конечное время --- это один шаг за один такт.

\begin{definition}[Элементарные шаг и такт]
\label{def:step-tick}
Пусть на~$\layerMicro$ заданы элементарный шаг решётки~$\caStepSpace$,
такт~$\caStepTime$ и тактовая скорость сигнала~$\caSignalSpeed$.
Говорят, что шаг и такт \emph{согласованы}, если
\begin{equation}
  \caSignalSpeed = \frac{\caStepSpace}{\caStepTime};
\end{equation}
тогда за один $\caStepTime$ сигнал не уходит дальше~$\caStepSpace$.
\textbf{Номер такта} --- это $\symTau\in\spaceZ$;
после планковской идентификации (\cref{def:si-bridge}) $t_{\mathrm{SI}}(\symTau)=\symTau\,\caStepTime$.
\end{definition}

\begin{definition}[Окрестность одного такта]
\label{def:neighborhood}
Пусть $x\in\carrierLattice$.
\textbf{Окрестность} узла~$x$ на один такт --- это множество узлов, до которых из~$x$
добираются ровно за один такт:
\[
  N(x)\subset\carrierLattice.
\]
\end{definition}

Требования спина и сохранения нормы уже записаны в квантовой механике.
На узле они остаются теми же:
значение --- спинор из двух компонент, норма --- сумма квадратов модулей.

\begin{definition}[Поле и норма]
\label{def:field-a3}
Пусть на каждом узле $x\in\carrierLattice$ в такт~$t$ задан спинор
\[
  z(x,t)=(z_1,z_2)^\top\in\spaceC^2.
\]
Один такой узел и есть \textbf{элементарная ячейка} $\caPlanckCell$.
\textbf{Плотность поля} на узле --- $|z_1|^2+|z_2|^2$;
\textbf{фазовые степени свободы} --- аргументы $z_k/|z_k|$, которые эволюция перемешивает.
\textbf{Глобальная конфигурация} в такт~$t$ записывается как
\[
  \symCapitalPsi^t=\{z(x,t)\}_{x\in\carrierLattice}.
\]
\textbf{Суммарную норму} (аналог $\langle\wavePsi|\wavePsi\rangle$) определяют формулой
\[
  N(t):=\sum_{x\in\carrierLattice} z^\dagger(x,t)\,z(x,t)=\sum_x\bigl(|z_1|^2+|z_2|^2\bigr).
\]
\end{definition}

Дальше --- не новый «язык'', а те же требования, записанные без континуума:
дискретная фаза, конечный регистр на узле, обратимый локальный шаг.
Геометрия~\latGCK\ $\carrierLattice$ и конкретный закон~$g$ выбираются в следующих главах;
здесь фиксируется, \emph{что} осталось после QED и планковского счёта.

\section{Предшествующая линия: от «вычисляющего пространства» к локальному автомату}
\label{sec:thooft-programme}

Мысль о мире как вычислении не наша и не его. В линии, которую цитирует эта книга
(приложение~\ref{ch:literature}), первым стоит фон Нейман (\textcite{vonneumann1966}):
самовоспроизводящиеся автоматы как модель того, что природа может делать без внешнего
конструктора. Дальше --- Цузе (\textcite{zuse1969}, \emph{Rechnender Raum}): физика как
дискретная процедура на сетке, время как счётчик тактов. Между ними и квантовой механикой
лежит линия обратимых вычислений --- Ландауэр, Беннет, Фредкин, Тоффоли, --- без которой
обратимый шаг не был бы сформулирован вовсе. На этом фундаменте Г.~’т~Хоофт
(\textcite{tHooft2016}) поставил программу, к тем же выводам которой эта книга пришла
**независимо**: квантовую механику можно и нужно строить снизу --- онтологические состояния
как перестановка на конечном множестве без суперпозиций, шаблонные состояния как линейные
статистические конструкции наблюдателя, а видимая нелокальность как следствие неполноты
шаблона, а не сигнала быстрее света.
В его терминологии~$\Mlayer$ --- онтологический слой, $\Tlayer$ --- шаблонный.

Приоритет ни в одном из этих звеньев не заявляется; полная линия с ссылками ---
в приложении~\ref{ch:literature}. Наше здесь одно звено --- носитель с законом шага.

Порядок здесь существен и его стоит зафиксировать: конструкция (\cref{ch:descent}--\cref{ch:si-sm})
построена раньше, чем эта монография была прочитана, --- расхождение в восемь недель по
журналу разработки. Поэтому совпадение выводов работает как \emph{независимый путь},
а не как пересказ: заранее объявленная реализация чужой программы ничего не доказывает,
а схождение, замеченное после постройки, --- доказывает.

Пять открытых проблем, перечисленных им самим (выбор решётки и модуля кольца;
разгон мод на краю зоны Бриллюэна; расходимость ряда Бейкера--Кэмпбелла--Хаусдорфа;
классы информационной эквивалентности и стрела времени; конструктивная реализация
Стандартной модели), разобраны в томе «Вычисления».
Точек, где независимый путь расходится с его программой, ровно одна --- и она решающая:
носитель.

Ошибка была в решётке. На декартовой сетке $\mathbb{Z}^d$ дисперсия
$4\sin^2(\omega/2)=\sum_i 2(1-\cos k_i)$ при $d>1$ на краю зоны Бриллюэна
превышает $4$ и переводит $\cos\omega$ в гиперболический косинус: моды разгоняются
экспоненциально (это его собственный \S\,17.2.1).
Плотнейшая изотропная упаковка с двенадцатью светоподобными связями
(кубооктаэдр, гл.~\ref{ch:carrier}) снимает разгон: вес связи $1/12$,
фактор вписанной сферы $\caKappaGeom^2=1/2$ и насыщающий импульс с порогом
Гейзенберга удерживают поворот на единичной окружности во всей трёхмерной зоне.

